\documentclass[aspectratio=169, 10pt]{beamer}

% ============================================================================
% PACKAGES & ENCODING (Configured for pdflatex + Lean 4 Unicode)
% ============================================================================
\usepackage[utf8]{inputenc}
\usepackage[T1]{fontenc}
\usepackage{lmodern}
\usepackage{amsmath, amssymb, amsfonts, mathtools}
\usepackage{booktabs}
\usepackage{xcolor}
\usepackage{listings}

% Global Unicode mappings for normal text and \texttt{...}
\DeclareUnicodeCharacter{21A6}{\ensuremath{\mapsto}}
\DeclareUnicodeCharacter{2192}{\ensuremath{\to}}
\DeclareUnicodeCharacter{2190}{\ensuremath{\leftarrow}}
\DeclareUnicodeCharacter{2194}{\ensuremath{\leftrightarrow}}
\DeclareUnicodeCharacter{211D}{\ensuremath{\mathbb{R}}}
\DeclareUnicodeCharacter{2102}{\ensuremath{\mathbb{C}}}
\DeclareUnicodeCharacter{2115}{\ensuremath{\mathbb{N}}}
\DeclareUnicodeCharacter{211A}{\ensuremath{\mathbb{Q}}}
\DeclareUnicodeCharacter{210B}{\ensuremath{\mathcal{H}}}
\DeclareUnicodeCharacter{2022}{\ensuremath{\bullet}}
\DeclareUnicodeCharacter{2200}{\ensuremath{\forall}}
\DeclareUnicodeCharacter{2203}{\ensuremath{\exists}}
\DeclareUnicodeCharacter{2227}{\ensuremath{\land}}
\DeclareUnicodeCharacter{2228}{\ensuremath{\lor}}
\DeclareUnicodeCharacter{00AC}{\ensuremath{\neg}}
\DeclareUnicodeCharacter{2208}{\ensuremath{\in}}
\DeclareUnicodeCharacter{2260}{\ensuremath{\neq}}
\DeclareUnicodeCharacter{2264}{\ensuremath{\le}}
\DeclareUnicodeCharacter{2211}{\ensuremath{\sum}}
\DeclareUnicodeCharacter{221A}{\ensuremath{\sqrt{}}}
\DeclareUnicodeCharacter{207B}{\ensuremath{^-}}
\DeclareUnicodeCharacter{00B9}{\ensuremath{^1}}
\DeclareUnicodeCharacter{00B2}{\ensuremath{^2}}
\DeclareUnicodeCharacter{2113}{\ensuremath{\ell}}
\DeclareUnicodeCharacter{2295}{\ensuremath{\oplus}}
\DeclareUnicodeCharacter{00D7}{\ensuremath{\times}}
\DeclareUnicodeCharacter{03B1}{\ensuremath{\alpha}}
\DeclareUnicodeCharacter{03B2}{\ensuremath{\beta}}
\DeclareUnicodeCharacter{03BC}{\ensuremath{\mu}}
\DeclareUnicodeCharacter{03C8}{\ensuremath{\psi}}
\DeclareUnicodeCharacter{03A8}{\ensuremath{\Psi}}
\DeclareUnicodeCharacter{2080}{\ensuremath{_0}}
\DeclareUnicodeCharacter{2081}{\ensuremath{_1}}
\DeclareUnicodeCharacter{2082}{\ensuremath{_2}}
\DeclareUnicodeCharacter{1D65}{\ensuremath{_v}}
\DeclareUnicodeCharacter{2016}{\ensuremath{\|}}
\DeclareUnicodeCharacter{00B7}{\ensuremath{\cdot}}
\DeclareUnicodeCharacter{22A2}{\ensuremath{\vdash}}
\DeclareUnicodeCharacter{00B7}{\ensuremath{\cdot}}
\DeclareUnicodeCharacter{208A}{\ensuremath{_+}}
\DeclareUnicodeCharacter{208B}{\ensuremath{_-}}
\DeclareUnicodeCharacter{03B5}{\ensuremath{\varepsilon}} % ε
\DeclareUnicodeCharacter{03BB}{\ensuremath{\lambda}}     % λ
\DeclareUnicodeCharacter{03C9}{\ensuremath{\omega}}      % ω
\DeclareUnicodeCharacter{03D5}{\ensuremath{\phi}}        % ϕ
\DeclareUnicodeCharacter{03C6}{\ensuremath{\varphi}}     % φ
\DeclareUnicodeCharacter{2208}{\ensuremath{\in}}         % ∈
\DeclareUnicodeCharacter{2211}{\ensuremath{\sum}}        % ∑
\DeclareUnicodeCharacter{27E8}{\ensuremath{\langle}}     % ⟨
\DeclareUnicodeCharacter{27E9}{\ensuremath{\rangle}}     % ⟩
\DeclareUnicodeCharacter{221A}{\ensuremath{\sqrt{}}}
\DeclareUnicodeCharacter{1D65}{\ensuremath{_v}}
\DeclareUnicodeCharacter{207B}{\ensuremath{^-}}
\DeclareUnicodeCharacter{00B9}{\ensuremath{^1}}
\DeclareUnicodeCharacter{2022}{\ensuremath{\bullet}}

% ============================================================================
% BEAMER THEME & COLORS
% ============================================================================
\usetheme{Madrid}
\usecolortheme{default}
\setbeamertemplate{navigation symbols}{}

\definecolor{leanblue}{RGB}{0, 90, 180}
\definecolor{leankeyword}{RGB}{150, 0, 130}
\definecolor{leancomment}{RGB}{90, 130, 90}
\definecolor{leanbg}{RGB}{248, 249, 252}

% ============================================================================
% LEAN 4 LISTINGS CONFIGURATION
% ============================================================================
\lstdefinelanguage{Lean4}{
  extendedchars=false,
  columns=fullflexible,
  keepspaces=true,
  upquote=true,
  basicstyle=\ttfamily\footnotesize,
  backgroundcolor=\color{leanbg},
  frame=single,
  rulecolor=\color{black!15},
  xleftmargin=4pt,
  xrightmargin=4pt,
  aboveskip=6pt,
  belowskip=4pt,
  breaklines=true,
  showstringspaces=false,
  morekeywords=[1]{
    def, theorem, lemma, example, inductive, structure, class, instance,
    where, by, fun, match, with, if, then, else, let, have, show, from,
    open, scoped, variable, import, noncomputable, omit, private, local
  },
  morekeywords=[2]{
    Prop, Type, Nat, Int, Real, Complex, Fin, List, Or, And, Eq, False, True,
    Matrix, EuclideanSpace, Orthonormal, OrthonormalBasis, Memlp, rootsOfUnity
  },
  morekeywords=[3]{
    intro, exact, apply, rw, simp, simp_rw, unfold, rcases, cases,
    constructor, ext, ring, norm_num, norm_cast, field_simp, obtain, use, suffices, simpa
  },
  keywordstyle=[1]\color{leanblue}\bfseries,
  keywordstyle=[2]\color{teal!80!black}\bfseries,
  keywordstyle=[3]\color{leankeyword}\bfseries,
  comment=[l]{--},
  morecomment=[s]{/-}{-/},
  commentstyle=\color{leancomment}\itshape,
  literate=
    {ℂ}{{$\mathbb{C}$}}1
    {ℝ}{{$\mathbb{R}$}}1
    {ℚ}{{$\mathbb{Q}$}}1
    {ℕ}{{$\mathbb{N}$}}1
    {ℋ}{{$\mathcal{H}$}}1
    {ℓ}{{$\ell$}}1
    {Ψ}{{$\Psi$}}1
    {ψ}{{$\psi$}}1
    {μ}{{$\mu$}}1
    {α}{{$\alpha$}}1
    {β}{{$\beta$}}1
    {₀}{{$_0$}}1
    {₁}{{$_1$}}1
    {₂}{{$_2$}}1
    {ᵥ}{{$_v$}}1
    {⁻¹}{{$^{-1}$}}2
    {²}{{$^2$}}1
    {→}{{$\to$}}1
    {←}{{$\leftarrow$}}1
    {↔}{{$\leftrightarrow$}}1
    {=>}{{$\mapsto$}}2
    {∀}{{$\forall$}}1
    {∃}{{$\exists$}}1
    {∈}{{$\in$}}1
    {≠}{{$\neq$}}1
    {≤}{{$\le$}}1
    {•}{{$\bullet$}}1
    {∧}{{$\land$}}1
    {∨}{{$\lor$}}1
    {¬}{{$\neg$}}1
    {×}{{$\times$}}1
    {⊕}{{$\oplus$}}1
    {∑}{{$\sum$}}1
    {√}{{$\sqrt{}$}}1
    {‖}{{$\|$}}1
    {·}{{$\cdot$}}1
    {⊢}{{$\vdash$}}1
    {₊}{{$\ensuremath{_{+}}$}}1
    {₋}{{$\ensuremath{_{-}}$}}1
    {ε}{{$\varepsilon$}}1
    {λ}{{$\lambda$}}1
    {ω}{{$\omega$}}1
    {∈}{{$\in$}}1
    {∑}{{$\sum$}}1
    {⟨}{{$\langle$}}1
    {⟩}{{$\rangle$}}1
    {√}{{$\sqrt{}$}}1
    {ᵥ}{{$_v$}}1
    {⁻}{{$^-$}}1
    {¹}{{$^1$}}1
    {•}{{$\bullet$}}1
}

% ============================================================================
% TITLE METADATA (SLIDE 1)
% ============================================================================
\title[Lean 4 for Physicists]{Formalizing Physics in Lean 4}
\author[M. Pacholski]{Micha\l{} Pacholski}
\institute[Google / FUW Alumnus]{
  Software Engineer, Google Warsaw \\
  \smallskip
  \textit{Faculty of Physics, University of Warsaw (FUW)}
}
\date{30 October 2026}

\begin{document}

% ----------------------------------------------------------------------------
% SLIDE 1: Title Slide
% ----------------------------------------------------------------------------
\begin{frame}
  \titlepage
\end{frame}

% ============================================================================
% SECTION 1: INTRODUCTION & MOTIVATION
% ============================================================================
\section{Introduction \& Motivation}

% ----------------------------------------------------------------------------
% SLIDE 2: Self-Introduction
% ----------------------------------------------------------------------------
\begin{frame}{About Me: From Condensed Matter to Software \& Back}
  \textbf{Academic \& Industry Trajectory}
  \begin{itemize}
    \item \textbf{University of Warsaw (FUW):} BSc \& MSc in Physics.
    \item \textbf{Leiden University:} PhD in Theoretical Condensed Matter Physics.
    \item \textbf{MPI-PKS Dresden:} Postdoctoral Researcher in Condensed Matter Theory.
    \item \textbf{Google Warsaw (2024--present):} Software Engineer.
  \end{itemize}

  \vspace{0.3cm}
  \begin{block}{Exploring Lean 4}
    \begin{itemize}
      \item Around 2024, I began exploring \textbf{Lean 4} as a personal passion project.
      \item It sits right at the intersection of my two worlds: theoretical physics and modern software engineering.
    \end{itemize}
  \end{block}
\end{frame}

% ----------------------------------------------------------------------------
% SLIDE 3: Presentation Outline
% ----------------------------------------------------------------------------
\begin{frame}{Outline of Today's Seminar}
  \begin{enumerate}
    \setlength{\itemsep}{6pt}
    \item \textbf{Why Learn Lean 4?} \\
      {\small Verification in the age of AI and importing mathematics like a software codebase.}

    \item \textbf{Fundamentals of Functional Programming} \\
      {\small Pure expressions, immutability, curried functions, and higher-order operators.}

    \item \textbf{Basic Lean 4 Syntax \& Data Types} \\
      {\small \texttt{\#check}, \texttt{\#eval}, \texttt{def}, products, sums, structures, and inductive types (\texttt{Nat}, \texttt{List}).}

    \item \textbf{The Curry--Howard Isomorphism} \\
      {\small Propositions as types, proofs as functional programs, and dependent types.}

    \item \textbf{Proving Theorems in Lean: Term Mode vs.\ Tactics} \\
      {\small Direct term construction vs.\ interactive tactic proofs (\texttt{intro}, \texttt{rw}, \texttt{simp}).}

    \item \textbf{Condensed Matter Case Study: The 1D Tight-Binding Chain} \\
      {\small Periodic boundary conditions via \texttt{Fin N}, Bloch waves via \texttt{rootsOfUnity}, and an exact orthonormal eigenbasis.}
  \end{enumerate}
\end{frame}

% ----------------------------------------------------------------------------
% SLIDE 4: Why Learn Lean 4? (Part I — Verification in the Age of AI)
% ----------------------------------------------------------------------------
\begin{frame}[fragile]{Why Learn Lean 4? I. Verification in the Age of AI}
  \begin{columns}[T]
    \begin{column}{0.50\textwidth}
      \textbf{Informal Derivations Can Be Hallucinated}
      \begin{itemize}
        \item LLMs can generate natural-language mathematical arguments and informal derivations that read smoothly to a human.
        \item Yet the entire argument---or a critical link in the reasoning---can be completely hallucinated.
        \item Checking an LLM's informal derivation line-by-line is slow, exhausting, and error-prone.
      \end{itemize}
    \end{column}

    \begin{column}{0.46\textwidth}
      \begin{alertblock}{The Solution: Dual Output}
        \begin{itemize}
          \item Require the AI to produce \textbf{both} an informal explanation \emph{and} formal \textbf{Lean 4 code}.
          \item Lean's kernel automatically verifies that the proof is logically sound.
          \item \textbf{Your job as a physicist:} Check that the Lean \texttt{def} and \texttt{theorem} statements actually match the informal claim!
        \end{itemize}
      \end{alertblock}
    \end{column}
  \end{columns}

  \vspace{0.2cm}
\begin{lstlisting}[language=Lean4]
-- You don't need to read the proof behind "by"; just verify the statement matches the claim:
theorem isEigenvector_planeWave :
    (hamiltonian N μ t) *ᵥ planeWave N u =
    (μ + t * u.val⁻¹ + star t * u.val) • planeWave N u := by ...
\end{lstlisting}
\end{frame}


% ----------------------------------------------------------------------------
% SLIDE 5: Why Learn Lean 4? (Part II — Physics as an Importable Codebase)
% ----------------------------------------------------------------------------
\begin{frame}[fragile]{Why Learn Lean 4? II. Physics as an Importable Codebase}
  \begin{columns}[T]
    \begin{column}{0.48\textwidth}
      \textbf{How Software Scales (e.g.\ at Google)}
      \begin{itemize}
        \item Thousands of engineers collaborate on one massive codebase over many years.
        \item You don't rewrite basic utilities---you \texttt{import} existing libraries, and the compiler guarantees compatibility.
      \end{itemize}
    \end{column}

    \begin{column}{0.48\textwidth}
      \textbf{How Theoretical Physics Works Today}
      \begin{itemize}
        \item Results live in static PDF papers with varying notations and implicit assumptions.
        \item Every new paper or textbook re-derives the same sums, integrals, and linear algebra from scratch.
      \end{itemize}
    \end{column}
  \end{columns}

  \vspace{0.15cm}
  \textbf{Lean's \texttt{Mathlib}:} A unified codebase of pure mathematics (linear algebra, analysis, algebra) that we can directly \texttt{import} to build a shared codebase for physics:

\begin{lstlisting}[language=Lean4]
import Mathlib.LinearAlgebra.Matrix.Hermitian
import Mathlib.RingTheory.RootsOfUnity.Complex
import Mathlib.Analysis.InnerProductSpace.PiL2
\end{lstlisting}
\end{frame}


% ============================================================================
% SECTION 2: FUNDAMENTALS OF FUNCTIONAL PROGRAMMING
% ============================================================================
\section{Fundamentals of Functional Programming}

% ----------------------------------------------------------------------------
% SLIDE 6: Fundamentals of Functional Programming
% ----------------------------------------------------------------------------
\begin{frame}[fragile]{Fundamentals of Functional Programming}
  \begin{columns}[T]
    \begin{column}{0.48\textwidth}
      \textbf{Imperative (Python, C++, Fortran)}
      \begin{itemize}
        \item Sequence of \textbf{state mutations} (\texttt{x = x + 1}) and in-place side effects.
      \end{itemize}
    \end{column}

    \begin{column}{0.48\textwidth}
      \textbf{Functional (Lean 4, Haskell, OCaml)}
      \begin{itemize}
        \item Evaluation of \textbf{pure, immutable expressions}: \texttt{f x} always yields the same result.
      \end{itemize}
    \end{column}
  \end{columns}

  \vspace{0.15cm}
  \textbf{Key Functional Concepts in Lean 4:}
  \begin{itemize}
    \item \textbf{Everything has a type (\texttt{x : T}):} Numbers, functions, types, and even proofs.
    \item \textbf{First-class functions (\texttt{f : A → B}):} A function mapping inputs of type \texttt{A} to outputs of type \texttt{B} has type signature \texttt{A → B}. Functions are ordinary values and can be passed to higher-order functions (like operators in physics).
    \item \textbf{Currying (\texttt{A → B → C}):} Read as \texttt{A → (B → C)}---a two-argument function takes an \texttt{A} and returns a new function of type \texttt{B → C}.
  \end{itemize}

\begin{lstlisting}[language=Lean4]
-- Anonymous function (lambda) of type ℕ → ℕ:
fun (x : ℕ) => x + 1
-- Higher-order function of type (ℕ → ℕ) → ℕ → ℕ:
def applyTwice (f : ℕ → ℕ) (x : ℕ) : ℕ := f (f x)
\end{lstlisting}
\end{frame}

% ============================================================================
% SECTION 3: BASIC LEAN 4 SYNTAX & DATA TYPES
% ============================================================================
\section{Basic Lean 4 Syntax \& Data Types}

% ----------------------------------------------------------------------------
% SLIDE 7: Basic Lean Syntax — def, #check, and #eval
% ----------------------------------------------------------------------------
\begin{frame}[fragile]{Basic Lean 4 Syntax: \texttt{def}, \texttt{\#check}, and \texttt{\#eval}}
  \textbf{Interacting with the Lean 4 Compiler:}
  \begin{itemize}
    \item \texttt{def} binds a name to a value or function with explicit parameter and return types.
    \item \texttt{\#check} asks Lean to infer and display the \textbf{type signature} of an expression without evaluating it.
    \item \texttt{\#eval} executes a computable expression and prints its resulting \textbf{value}.
  \end{itemize}

\begin{lstlisting}[language=Lean4]
-- Define a constant and a curried two-argument function:
def n : ℕ := 5
def add (x : ℕ) (y : ℕ) : ℕ := x + y

-- Inspect types with #check (notice how partial application works!):
#check n          -- n : ℕ
#check add        -- add : ℕ → ℕ → ℕ
#check add 2      -- add 2 : ℕ → ℕ   (partially applied function!)
#check add 2 3    -- add 2 3 : ℕ

-- Compute values with #eval:
#eval add 2 3     -- 5
\end{lstlisting}
\end{frame}

% ----------------------------------------------------------------------------
% SLIDE 8: Combining Data I — Product Types and Structures
% ----------------------------------------------------------------------------
\begin{frame}[fragile]{Combining Data I: Product Types and \texttt{structure}}
  \textbf{Holding Multiple Values Simultaneously:}
  \begin{itemize}
    \item \textbf{Product Types (\texttt{A × B}):} Hold \textbf{both} an \texttt{A} \emph{and} a \texttt{B} as an anonymous pair \texttt{(a, b)}. Fields are accessed via \texttt{.1} and \texttt{.2} (or \texttt{.fst} and \texttt{.snd}).
    \item \textbf{Structures (\texttt{structure}):} Define a custom, named product type with readable field names instead of positional indices.
  \end{itemize}

  \vspace{0.15cm}
\begin{lstlisting}[language=Lean4]
-- Anonymous product type (pair):
def p : ℕ × ℝ := (3, 2.5)
#eval p.1    -- 3
#eval p.2    -- 2.5

-- Named structure:
structure Point2D where
  x : ℝ
  y : ℝ

def pt : Point2D := { x := 1.0, y := -1.0 }
#check pt.x  -- pt.x : ℝ
\end{lstlisting}
\end{frame}

% ----------------------------------------------------------------------------
% SLIDE 9: Combining Data II — Sum Types and Pattern Matching
% ----------------------------------------------------------------------------
\begin{frame}[fragile]{Combining Data II: Sum Types (\texttt{A ⊕ B})}
  \textbf{Representing Alternatives (Disjoint Unions):}
  \begin{itemize}
    \item \textbf{Sum Types (\texttt{A ⊕ B} or \texttt{Sum A B}):} Hold \textbf{either} a value of type \texttt{A} (tagged with \texttt{Sum.inl a}) \emph{or} a value of type \texttt{B} (tagged with \texttt{Sum.inr b}).
    \item \textbf{Pattern Matching (\texttt{match ... with}):} To use a sum type, we must provide a branch for every possible constructor---the compiler guarantees no case is forgotten!
  \end{itemize}

  \vspace{0.15cm}
\begin{lstlisting}[language=Lean4]
-- Constructing values of a Sum type (ℕ ⊕ ℝ):
def s1 : ℕ ⊕ ℝ := Sum.inl 42
def s2 : ℕ ⊕ ℝ := Sum.inr 3.14

-- Consuming a Sum type via pattern matching:
def toReal (s : ℕ ⊕ ℝ) : ℝ :=
  match s with
  | Sum.inl n => (n : ℝ)
  | Sum.inr r => r
\end{lstlisting}
\end{frame}

% ----------------------------------------------------------------------------
% SLIDE 10: Recursive Inductive Data Types — Nat and List
% ----------------------------------------------------------------------------
\begin{frame}[fragile]{Inductive Data Types: \texttt{Nat} and \texttt{List}}
  In Lean 4, even numbers and containers are not built-in ``magic'' primitives---they are defined from scratch as \textbf{recursive inductive types} (\texttt{inductive}):

  \vspace{0.1cm}
  \begin{columns}[T]
    \begin{column}{0.48\textwidth}
      \textbf{1. Peano Natural Numbers (\texttt{Nat} / \texttt{ℕ})}
      \begin{itemize}
        \item \texttt{zero}: the base number $0$.
        \item \texttt{succ n}: the successor $n + 1$ of an existing \texttt{n : Nat}.
        \item \texttt{2} is syntax sugar for \texttt{succ (succ zero)}.
      \end{itemize}
    \end{column}

    \begin{column}{0.48\textwidth}
      \textbf{2. Polymorphic Lists (\texttt{List α})}
      \begin{itemize}
        \item \texttt{nil}: the empty list \texttt{[]}.
        \item \texttt{cons head tail}: an element \texttt{head : α} prepended to \texttt{tail : List α} (written \texttt{head :: tail}).
      \end{itemize}
    \end{column}
  \end{columns}

  \vspace{0.1cm}
\begin{lstlisting}[language=Lean4]
-- How Nat and List are defined in Lean's standard library:
inductive Nat where
  | zero : Nat
  | succ (n : Nat) : Nat

inductive List (α : Type) where
  | nil  : List α
  | cons (head : α) (tail : List α) : List α
\end{lstlisting}
\end{frame}

% ----------------------------------------------------------------------------
% SLIDE 11: Pattern Matching & Recursion on Nat and List
% ----------------------------------------------------------------------------
\begin{frame}[fragile]{Pattern Matching \& Structural Recursion}
  Functions on inductive types are defined by \textbf{pattern matching} on each constructor. Lean ensures termination by checking that recursive calls only happen on smaller subterms:

  \vspace{0.15cm}
\begin{lstlisting}[language=Lean4]
-- Addition on Nat by recursion on the second argument:
def myAdd (m n : Nat) : Nat :=
  match n with
  | Nat.zero   => m
  | Nat.succ k => Nat.succ (myAdd m k)

-- Length of a polymorphic List (using [] for nil and x :: xs for cons):
def myLength (xs : List α) : Nat :=
  match xs with
  | []      => 0
  | _ :: tl => 1 + myLength tl

#eval myAdd 3 2             -- 5
#eval myLength [10, 20, 30] -- 3
\end{lstlisting}
\end{frame}

% ============================================================================
% SECTION 4: THE CURRY-HOWARD ISOMORPHISM
% ============================================================================
\section{The Curry--Howard Isomorphism}

% ----------------------------------------------------------------------------
% SLIDE 12: The Curry-Howard Isomorphism — Propositions as Types
% ----------------------------------------------------------------------------
\begin{frame}[fragile]{The Curry--Howard Isomorphism: Propositions as Types}
  In Lean 4, there is no boundary between \textbf{programming} and \textbf{proving}:
  \begin{itemize}
    \item A mathematical \textbf{proposition} $P$ is simply a \textbf{type} (\texttt{P : Prop}).
    \item A \textbf{proof} of $P$ is a \textbf{program (value)} $p$ of type $P$ (\texttt{p : P}).
    \item Checking if a proof is valid is literally \textbf{type-checking} the program!
  \end{itemize}

  \vspace{0.15cm}
  \begin{center}
  \small
  \begin{tabular}{lll}
    \textbf{Functional Programming} & \textbf{Lean 4 Syntax} & \textbf{Logic (Mathematics)} \\
    \hline
    Type & \texttt{P : Prop} & Proposition $P$ \\
    Value / Term of type $P$ & \texttt{h : P} & Proof of $P$ \\
    Empty / Unit type & \texttt{Empty}/\texttt{Unit} \uncover<2->{(\texttt{False}/\texttt{True})} & \uncover<2->{Falsehood / Truth} \\
    Function type & \texttt{P → Q} & \uncover<3->{Implication ($P \to Q$)} \\
    Product type & \texttt{A × B} \uncover<4->{(\texttt{P ∧ Q})} & \uncover<4->{Conjunction ($P \land Q$)} \\
    Sum type & \texttt{A ⊕ B} \uncover<5->{(\texttt{P ∨ Q})} & \uncover<5->{Disjunction ($P \lor Q$)}
  \end{tabular}
  \end{center}
\end{frame}

% ----------------------------------------------------------------------------
% SLIDE 13: Dependent Types — Quantifiers and Subtypes
% ----------------------------------------------------------------------------
\begin{frame}[fragile]{Dependent Types: Quantifiers ($\forall$, $\exists$) and Subtypes}
  In Lean 4, \textbf{types can depend on values}: a family of propositions $P(x)$ is simply a function \texttt{P : α → Prop} mapping each value \texttt{x : α} to a proposition \texttt{P x}.

  \vspace{0.15cm}
  \begin{itemize}
    \item \textbf{Universal Quantifier ($\forall x : \alpha,\, P(x)$) $\longleftrightarrow$ Dependent Function (\texttt{(x : α) → P x}):} \\
      A function that takes any input \texttt{x : α} and returns a proof of \texttt{P x} for that specific \texttt{x}.

    \item \textbf{Existential Quantifier ($\exists x : \alpha,\, P(x)$) $\longleftrightarrow$ Dependent Pair (\texttt{⟨x, h⟩}):} \\
      A structure packing a concrete witness \texttt{x : α} together with a proof \texttt{h : P x}.

    \item \textbf{Data + Proofs in Physics (e.g.\ \texttt{Fin N}):} \\
      A site on an $N$-site ring has type \texttt{Fin N}: a pair bundling a number \texttt{val : ℕ} with a proof \texttt{isLt : val < N}. Out-of-bounds indices cannot even be constructed!
  \end{itemize}

  \vspace{0.1cm}
\begin{lstlisting}[language=Lean4]
-- Universal quantifier (∀ n : ℕ, n = n) is a dependent function:
#check ∀ (n : ℕ), n = n     -- ∀ (n : ℕ), n = n : Prop

-- Fin N bundles a natural number with a proof that it is strictly less than N:
structure Fin (N : ℕ) where
  val  : ℕ
  isLt : val < N
\end{lstlisting}
\end{frame}

% ============================================================================
% SECTION 5: PROVING THEOREMS IN LEAN 4
% ============================================================================
\section{Proving Theorems in Lean 4}

% ----------------------------------------------------------------------------
% SLIDE 14: Proving by Direct Term Construction (Term Mode)
% ----------------------------------------------------------------------------
\begin{frame}[fragile]{Proving by Direct Term Construction (Term Mode)}
  Because a proof of $P$ is just a value of type $P$, we can prove theorems directly by writing a functional program of that type (\textbf{Term Mode}):
  \begin{itemize}
    \item \texttt{theorem name (args) : Statement := proof\_term} is identical to \texttt{def}, except its return type lives in \texttt{Prop}.
    \item \textbf{Implication ($P \to Q$):} Proven by a function \texttt{fun h => ...} mapping a proof \texttt{h : P} to a proof of \texttt{Q}.
    \item \textbf{Conjunction / Disjunction:} Proven using pair and sum constructors (\texttt{And.intro}, \texttt{Or.inl}).
  \end{itemize}

  \vspace{0.1cm}
\begin{lstlisting}[language=Lean4]
-- 1. Proving P → P using the identity function:
theorem id_proof (P : Prop) : P → P :=
  fun (h : P) => h

-- 2. Proving P ∧ Q → Q ∧ P by swapping a pair:
theorem and_swap (P Q : Prop) : P ∧ Q → Q ∧ P :=
  fun (h : P ∧ Q) => And.intro h.right h.left

-- 3. Proving P → P ∨ Q by injecting into the left branch:
theorem or_left (P Q : Prop) : P → P ∨ Q :=
  fun (h : P) => Or.inl h
\end{lstlisting}
\end{frame}

% ----------------------------------------------------------------------------
% SLIDE 15: Interactive Tactic Mode and the Lean Infoview
% ----------------------------------------------------------------------------
\begin{frame}[fragile]{Interactive Tactic Mode and the Lean Infoview}
  For complex theorems, writing raw proof terms by hand becomes unwieldy. Instead, we enter \textbf{Tactic Mode} using \texttt{by}:
  \begin{itemize}
    \item \textbf{Tactics} are metaprograms that build the underlying proof term step by step.
    \item At every cursor position, the \textbf{Lean Infoview} displays the current \textbf{Proof State}: local hypotheses at the top, and the remaining goal after the turnstile (\texttt{⊢}).
  \end{itemize}

  \vspace{0.1cm}
  \begin{columns}[T]
    \begin{column}{0.48\textwidth}
      \textbf{Tactic Script (\texttt{by}):}
\begin{lstlisting}[language=Lean4]
theorem and_swap (P Q : Prop) :
    P ∧ Q → Q ∧ P := by
  intro h
  constructor
  · exact h.right
  · exact h.left
\end{lstlisting}
    \end{column}

    \begin{column}{0.48\textwidth}
      \textbf{Lean Infoview (after \texttt{intro h}):}
\begin{lstlisting}[language=Lean4]
P Q : Prop
h : P ∧ Q
⊢ Q ∧ P

-- `constructor` splits ⊢ Q ∧ P
-- into two subgoals: ⊢ Q and ⊢ P
\end{lstlisting}
    \end{column}
  \end{columns}
\end{frame}

% ----------------------------------------------------------------------------
% SLIDE 16: Core Tactics — rw, simp, ring, and induction
% ----------------------------------------------------------------------------
\begin{frame}[fragile]{Core Tactics: \texttt{rw}, \texttt{simp}, \texttt{ring}, and \texttt{induction}}
  \begin{itemize}
    \item \texttt{rw [h]} (\textbf{Rewrite}): Substitutes equals for equals in the goal using equation \texttt{h} (use \texttt{rw [← h]} to rewrite right-to-left).
    \item \texttt{simp} (\textbf{Simplifier}): Automatically applies tagged simplification lemmas (\texttt{@[simp]}) to reduce the goal.
    \item \texttt{ring}: Automatically proves any polynomial identity in commutative (semi)rings.
    \item \texttt{induction n with}: Performs structural induction on an inductive type like \texttt{ℕ}.
  \end{itemize}

  \vspace{0.1cm}
\begin{lstlisting}[language=Lean4]
-- 1. Rewriting with a hypothesis h : a = b
example (a b c : ℝ) (h : a = b) : a + c = b + c := by
  rw [h]

-- 2. Automatic ring algebra (no manual expansion needed!)
example (x y : ℝ) : (x + y)^2 = x^2 + 2*x*y + y^2 := by
  ring

-- 3. Proving 0 + n = n by mathematical induction on n
theorem zero_add (n : ℕ) : 0 + n = n := by
  induction n with
  | zero      => rfl
  | succ k ih => rw [Nat.add_succ, ih]
\end{lstlisting}
\end{frame}

% ============================================================================
% SECTION 6: CASE STUDY — 1D TIGHT-BINDING CHAIN
% ============================================================================
\section{Case Study: 1D Tight-Binding Chain}

% ----------------------------------------------------------------------------
% SLIDE 17: 1D Tight-Binding Hamiltonian & Hermiticity
% ----------------------------------------------------------------------------
\begin{frame}[fragile]{1D Tight-Binding Chain: Hamiltonian \& Hermiticity}
  We model an $N$-site periodic ring (\texttt{Fin N}) with onsite chemical potential $\mu \in \mathbb{R}$ and nearest-neighbor hopping $t \in \mathbb{C}$ as an $N \times N$ matrix $H_{ij} = \mu\delta_{ij} + t\delta_{i,j+1} + t^*\delta_{j,i+1}$:

  \vspace{0.05cm}
\begin{lstlisting}[language=Lean4]
variable (N : ℕ) [NeZero N] (μ : ℝ) (t : ℂ)

def onsite (i j : Fin N) := if i = j then (μ : ℂ) else 0

def hopping_forward (i j : Fin N) := if i = j + 1 then t else 0

def hamiltonian_fun (i j : Fin N) :=
  onsite N μ i j + hopping_forward N t i j + star (hopping_forward N t j i)

def hamiltonian := Matrix.of (hamiltonian_fun N μ t)

theorem isHermitian_hamiltonian : (hamiltonian N μ t).IsHermitian := by
  ext i j
  unfold hamiltonian hamiltonian_fun
  simp only [Matrix.conjTranspose_apply, Matrix.of_apply, star_add,
    star_star, onsite_conjTranspose_eq_self, add_comm, add_assoc]
\end{lstlisting}
\end{frame}

% ----------------------------------------------------------------------------
% SLIDE 18: Bloch Plane Waves & Eigenvalue Theorem
% ----------------------------------------------------------------------------
\begin{frame}[fragile]{Bloch Plane Waves \& Eigenvalue Equation}
  Translational invariance is captured by indexing momenta with $N$-th roots of unity $u \in \mathbb{C}$ ($u^N = 1$, i.e.\ $u = e^{ik}$), defining normalized Bloch states $\psi_u(j) = u^j / \sqrt{N}$:

  \vspace{0.05cm}
\begin{lstlisting}[language=Lean4]
variable (u : rootsOfUnity N ℂ)

noncomputable def planeWaveFun (i : Fin N) : ℂ := ((u ^ i.val).val : ℂ) / √N

noncomputable def planeWave : EuclideanSpace ℂ (Fin N) where
  ofLp i := planeWaveFun N u i

theorem isEigenvector_planeWave :
    (hamiltonian N μ t) *ᵥ planeWave N u =
    (μ + t * u.val⁻¹ + star t * u.val) • planeWave N u := by
  ext x'
  rw [Matrix.mulVec_apply]
  unfold hamiltonian hamiltonian_fun planeWave planeWaveFun
         dotProduct hopping_forward onsite
  -- Simplify Kronecker deltas (Finset.sum_ite_eq) and modular exponents:
  ...
  · norm_num; norm_cast; simp [inv_u]; ring
\end{lstlisting}
\end{frame}

% ----------------------------------------------------------------------------
% SLIDE 19: Orthonormality & Complete Eigenbasis
% ----------------------------------------------------------------------------
\begin{frame}[fragile]{Orthonormality \& Complete Eigenbasis}
  Finally, we prove that the Bloch waves are orthonormal via the geometric series formula (\texttt{geom\_sum\_eq}) for $u' / u \neq 1$, and promote them to a complete \texttt{OrthonormalBasis} of $\mathbb{C}^N$:

  \vspace{0.05cm}
\begin{lstlisting}[language=Lean4]
theorem orthonormal_planeWaves : Orthonormal ℂ (planeWave N) := by
  constructor
  · intro u           -- Normalization: ∑ x, ‖planeWaveFun N u x‖^2 = 1
    ...
  · intro u u' h_ne   -- Orthogonality (u ≠ u'): geometric sum vanishes
    ...
    simp [Fin.sum_univ_eq_sum_range, geom_sum_eq h_ne']
    simp [div_pow]; norm_cast; simp [u_pow]

-- Since |rootsOfUnity N ℂ| = N = dim(ℂ^N), orthonormality implies completeness:
noncomputable def planeWaveBasis :
    Module.Basis (rootsOfUnity N ℂ) ℂ (EuclideanSpace ℂ (Fin N)) := by
  apply basisOfOrthonormalOfCardEqFinrank (orthonormal_planeWaves N)
  rw [Fintype.card_eq_nat_card, Complex.card_rootsOfUnity]; simp

noncomputable def planeWaveOrthonormalBasis :
    OrthonormalBasis (rootsOfUnity N ℂ) ℂ (EuclideanSpace ℂ (Fin N)) :=
  Basis.toOrthonormalBasis (planeWaveBasis N) (orthonormal_planeWavesBasis N)
\end{lstlisting}
\end{frame}

% ----------------------------------------------------------------------------
% SLIDE 20: Summary & Getting Started with Lean 4
% ----------------------------------------------------------------------------
\begin{frame}{Summary \& Resources for Getting Started}
  \textbf{Key Takeaways:}
  \begin{itemize}
    \item \textbf{Curry--Howard Isomorphism:} Propositions are types, proofs are programs, and proof verification is just type-checking.
    \item \textbf{Expressive Type System:} Dependent types (\texttt{Fin N}, \texttt{rootsOfUnity N ℂ}) encode physical constraints and symmetries directly in definitions.
    \item \textbf{Interactive Proving:} Tactics (\texttt{rw}, \texttt{simp}, \texttt{ring}) automate routine algebra while the Infoview tracks every logical assumption and goal.
  \end{itemize}

  \vspace{0.15cm}
  \textbf{Where to Start:}
  \begin{itemize}
    \item \textbf{Code from this Seminar:} \texttt{github.com/mpacholski/physics\_in\_lean}
    \item \textbf{The Natural Number Game:} Interactive browser puzzle game teaching tactics from scratch (\texttt{adam.math.hhu.de}).
    \item \textbf{Official Textbooks:} \emph{Theorem Proving in Lean 4}, \emph{Mathematics in Lean}, \& \emph{Functional Programming in Lean} (\texttt{lean-lang.org/documentation}).
    \item \textbf{Mathlib4 \& PhysLean (HepLean):} Community libraries formalizing mathematics and physics (\texttt{leanprover-community.github.io}).
    \item \textbf{Lean Zulip Chat:} Active, welcoming forum for questions and collaboration (\texttt{leanprover.zulipchat.com}).
  \end{itemize}
\end{frame}

\end{document}